\documentclass[12pt, a4paper]{article}

% ── PACKAGES ──────────────────────────────────────────────────
\usepackage[utf8]{inputenc}
\usepackage[T1]{fontenc}
\usepackage{geometry}
\geometry{margin=2.5cm}
\usepackage{hyperref}
\hypersetup{
    colorlinks=true,
    linkcolor=blue,
    urlcolor=blue,
    citecolor=blue,
    pdftitle={Entinostat Benefit in HR+ Breast Cancer
              Is a Luminal B-Specific Effect: The
              DNMT3A/HDAC2 Chromatin Lock Mechanism
              and Its Implications for NCT07235618},
    pdfauthor={Eric Robert Lawson},
    pdfkeywords={entinostat, HDACi, HDAC inhibitor,
                 Luminal B, LumB, LumA, breast cancer,
                 DNMT3A, HDAC2, chromatin lock,
                 TFF1, ESR1, ER output decoupling,
                 NCT07235618, Sun Yat-sen, PAM50,
                 subtype-specific, HR+, endocrine
                 therapy, CDK4/6 inhibitor, fulvestrant,
                 attractor geometry, Waddington landscape,
                 OrganismCore, GSE176078, METABRIC}
}
\usepackage{amsmath}
\usepackage{booktabs}
\usepackage{longtable}
\usepackage{array}
\usepackage{parskip}
\usepackage{microtype}
\usepackage{authblk}
\usepackage{titlesec}
\usepackage{fancyhdr}
\usepackage{xcolor}
\usepackage{float}
\usepackage{enumitem}
\usepackage{setspace}

% ── COLOURS ───────────────────────────────────────────────────
\definecolor{repobox}{RGB}{240,247,255}
\definecolor{findingbox}{RGB}{245,255,245}
\definecolor{convergentbox}{RGB}{255,250,235}
\definecolor{mechanismbox}{RGB}{250,245,255}
\definecolor{actionbox}{RGB}{230,255,230}
\definecolor{urgentbox}{RGB}{255,235,235}
\definecolor{novelbox}{RGB}{255,248,220}

% ── HEADER / FOOTER ───────────────────────────────────────────
\pagestyle{fancy}
\fancyhf{}
\lhead{\footnotesize CS-LIT-15: Entinostat ---
       LumB-Specific Benefit}
\rhead{\footnotesize OrganismCore \textbar{} 2026-03-06}
\rfoot{\small Page \thepage}
\renewcommand{\headrulewidth}{0.4pt}
\setlength{\headheight}{24pt}
\addtolength{\topmargin}{-10pt}

% ── SECTION FORMATTING ────────────────────────────────────────
\titleformat{\section}
  {\large\bfseries}{\thesection}{1em}{}[\titlerule]
\titleformat{\subsection}
  {\normalsize\bfseries}{\thesubsection}{1em}{}
\titleformat{\subsubsection}
  {\small\bfseries}{\thesubsubsection}{1em}{}

% ══════════════════════════════════════════════════════════════
% TITLE
% ══════════════════════════════════════════��═══════════════════
\title{
    \vspace{-1em}
    \textbf{Entinostat Benefit in HR$+$ Breast Cancer
    Is a Luminal~B-Specific Effect:}\\[0.5em]
    \large The DNMT3A/HDAC2 Chromatin Lock on ER
    Output as the Mechanistic Basis for Subtype-Enriched
    Patient Selection, with Implications for
    NCT07235618\\[0.4em]
    \normalsize \textit{OrganismCore Framework ---
    Document CS-LIT-15}\\[0.3em]
    \small Derived from Waddington Attractor Geometry
    of 19,542 Single Cancer Cells (GSE176078)\\[0.3em]
    \small DNMT3A/HDAC2 Coupling Replicated in METABRIC
    ($n=1{,}980$, $p=5.68\times10^{-56}$)
}

\author[1]{Eric Robert Lawson}
\affil[1]{
    OrganismCore\\
    \href{mailto:OrganismCore@proton.me}
         {OrganismCore@proton.me}\\
    ORCID:
    \href{https://orcid.org/0009-0002-0414-6544}
         {0009-0002-0414-6544}
}

\date{March 6, 2026}

% ══════════════════════════════════════════════════════════════
\begin{document}
% ══════════════════════════════════════════════════════════════

\maketitle
\thispagestyle{fancy}

% ── REPOSITORY POINTER ────────────────────────────────────────
\begin{center}
\colorbox{repobox}{
\begin{minipage}{0.88\textwidth}
\vspace{0.5em}
\centering
\textbf{Full computational record --- scripts, data
caches, reasoning artifacts, and raw logs:}\\[0.4em]
\href{https://github.com/Eric-Robert-Lawson/attractor-oncology}
{\texttt{https://github.com/Eric-Robert-Lawson/attractor-oncology}}
\\[0.3em]
\small Part of the OrganismCore BRCA Cross-Subtype
Analysis (BRCA-S8h, 2026-03-05).\\
All predictions locked from attractor geometry
\textbf{before} any literature review was performed.
\vspace{0.5em}
\end{minipage}}
\end{center}

\vspace{0.8em}

% ── URGENCY BOX ───────────────────────────────────────────────
\begin{center}
\colorbox{urgentbox}{
\begin{minipage}{0.88\textwidth}
\vspace{0.5em}
\centering
\textbf{TIME-SENSITIVE NOTE}\\[0.4em]
\small
NCT07235618 --- Phase~II entinostat $+$ fulvestrant
in HR$+$/HER2$-$ breast cancer post-CDK4/6i failure,
Sun Yat-sen University Cancer Center, Guangzhou.
Start date: January 2026. $N=50$. Primary endpoint:
PFS. Currently recruiting.\\[0.3em]
This trial is a de facto Luminal~B-enriched
population. Its sole biomarker endpoint is PBMC
histone acetylation. It does not currently include
LumA vs LumB PAM50 stratification or TFF1/ESR1
IHC as a correlative endpoint.\\[0.3em]
\textbf{The LumB subtype stratification described
in this document could be added as a correlative
endpoint on already-planned tissue without
protocol amendment, change to primary endpoint,
or additional patient burden. This requires
PAM50 subtyping or FOXA1/EZH2 IHC ratio
on baseline core biopsy tissue.}
\vspace{0.5em}
\end{minipage}}
\end{center}

\vspace{0.8em}

% ── CONVERGENT BOX ────────────────────────────────────────────
\begin{center}
\colorbox{convergentbox}{
\begin{minipage}{0.88\textwidth}
\vspace{0.5em}
\textbf{CONVERGENT EVIDENCE:}\\[0.3em]
Entinostat is now \textbf{approved in China}
(NMPA, April 2024) for HR$+$/HER2$-$ advanced
breast cancer after prior endocrine therapy ---
the first regulatory approval for an HDACi in
breast cancer. A May 2025 meta-analysis
(Springer, $n=1{,}371$) confirmed PFS improvement
(HR$=0.80$, $p=0.01$) with no OS benefit.\\[0.3em]
The drug is clinically active in the HR$+$ class.
The framework's novel contribution is the
mechanistic explanation of \textit{why} the effect
is modest and who benefits most: the benefit is
LumB-specific, driven by the DNMT3A/HDAC2
chromatin lock on ER output that is present in
LumB and absent in LumA. Treating HR$+$ as a
class dilutes the LumB signal with non-responding
LumA patients.
\vspace{0.5em}
\end{minipage}}
\end{center}

\vspace{0.8em}

% ══════════════════════════════════════════════════════════════
\begin{abstract}
% ══════════════════════════════════════════════════════════════
\noindent
Entinostat is now approved in China (NMPA,
April 2024) for HR$+$/HER2$-$ advanced breast
cancer, and a May 2025 meta-analysis of four
RCTs ($n=1{,}371$) confirmed PFS benefit
(HR$=0.80$, $p=0.01$) with no OS benefit. No
published trial has stratified entinostat
benefit by Luminal~A (LumA) versus Luminal~B
(LumB) PAM50 subtype. This document reports
the OrganismCore attractor geometry finding that
entinostat benefit is a LumB-specific phenomenon,
not a class effect across all HR$+$ disease.

The mechanistic basis, derived from Waddington
landscape analysis of 19,542 single cancer cells
(GSE176078) and replicated in METABRIC
($n=1{,}980$), is a LumB-specific co-repressor
circuit: DNMT3A and HDAC2 are co-expressed in
LumB at $r=+0.267$, compared to $r=+0.071$ in
LumA ($p=5.68\times10^{-56}$ for the difference).
This circuit acts as a chromatin lock on ER
output, suppressing downstream oestrogen-responsive
gene transcription (TFF1) despite elevated
ESR1 transcript --- the TFF1/ESR1 decoupling
replicated in METABRIC ($p=0.0019$). Entinostat
(class~I HDAC inhibitor) disrupts this circuit
specifically in LumB because the circuit is
present in LumB and absent in LumA. LumA patients
do not have the lock; they have no mechanism
through which entinostat improves ER output. The
modest HR$=0.80$ effect in the class-level
meta-analysis is consistent with a strong LumB
signal diluted by non-responding LumA patients.

The prediction is directly testable in
NCT07235618 (entinostat $+$ fulvestrant,
post-CDK4/6i failure, Sun Yat-sen University,
$N=50$, start January 2026), a trial whose
post-CDK4/6i design selects a de facto
LumB-enriched population. PAM50 subtype
stratification or FOXA1/EZH2 IHC ratio
(CS-LIT-22,
\href{https://doi.org/10.5281/zenodo.18892788}
{DOI: 10.5281/zenodo.18892788}) could be added
as a correlative endpoint on baseline tissue
without protocol amendment.

This subtype-specific mechanistic prediction ---
LumB specifically, DNMT3A/HDAC2 mechanism,
not a class effect --- has no direct published
equivalent as of 2026-03-06. Verdict:
CONVERGENT-NOVEL.
\end{abstract}

\tableofcontents
\newpage

% ══════════════════════════════════════════════════════════════
\section{Context and Origin}
% ══════════════════════════════════════════════════════════════

This document covers CS-LIT-15 from the
OrganismCore BRCA Cross-Subtype Literature Check
(BRCA-S8h, 2026-03-05). The entinostat
prediction was derived from attractor geometry
before any literature review was performed.
The literature check found that entinostat is
now clinically approved, confirming the drug is
active in the target population (HR$+$ breast
cancer). The subtype-specific prediction remains
novel and untested.

\subsection{Position in the Published Series}

\begin{itemize}
    \item \textbf{CS-LIT-9/23 (TFF1/ESR1
    Decoupling as LumB HDACi Biomarker)}\\
    \href{https://doi.org/10.5281/zenodo.18884234}
    {DOI: 10.5281/zenodo.18884234}\\
    \textit{The downstream consequence of the
    DNMT3A/HDAC2 lock --- ER output suppressed
    in LumB. CS-LIT-15 is the drug intervention
    document for CS-LIT-9/23.}

    \item \textbf{CS-LIT-1 (FOXA1/EZH2 Ratio
    Ordering Axis)}\\
    \href{https://doi.org/10.5281/zenodo.18883922}
    {DOI: 10.5281/zenodo.18883922}\\
    \textit{The continuous ordering that places
    LumB between LumA and HER2, with a distinct
    chromatin mechanism from LumA.}

    \item \textbf{CS-LIT-2 (Six Lock Type
    Classification)}\\
    \href{https://doi.org/10.5281/zenodo.18884158}
    {DOI: 10.5281/zenodo.18884158}\\
    \textit{LumB is Lock Type~1 with an epigenetic
    brake (DNMT3A/HDAC2). Entinostat addresses
    the brake. CS-LIT-15 is the drug evidence
    document for this lock type.}

    \item \textbf{CS-LIT-22 (FOXA1/EZH2 Dual
    IHC as Point-of-Care Tool)}\\
    \href{https://doi.org/10.5281/zenodo.18892788}
    {DOI: 10.5281/zenodo.18892788}\\
    \textit{The FOXA1/EZH2 IHC ratio classifies
    LumA vs LumB at the pathology lab level ---
    the patient selection tool for NCT07235618
    stratification.}
\end{itemize}

% ══════════════════════════════════════════════════════════════
\section{The LumA/LumB Paradox: More ER,
Less Sensitivity}
% ══════════════════════════════════════════════════════════════

Luminal~B (LumB) breast cancer presents a
well-recognised paradox: it has higher ESR1
(oestrogen receptor) transcript than Luminal~A
(LumA), yet is more resistant to endocrine
therapy and has worse prognosis. No published
biomarker framework has provided a causal
mechanistic explanation for this paradox.

The attractor geometry framework provides the
explanation: in LumB, the ER gene is present at
higher transcript levels, but its transcriptional
output is suppressed by a co-repressor circuit
involving DNMT3A and HDAC2. The receptor is there.
The output is locked. This is ER output decoupling.
It is not ER absence. It is ER silencing.

The specific prediction that follows: HDAC
inhibition --- which dissolves HDAC2-dependent
co-repressor activity --- restores ER output
specifically in LumB, because the HDAC lock is
present in LumB and absent in LumA.

% ══════════════════════════════════════════════════════════════
\section{The DNMT3A/HDAC2 Chromatin Lock}
% ══════════════════════════════════════════════════════════════

\subsection{The Co-Expression Coupling}

From Waddington attractor geometry of 19,542
single cancer cells (GSE176078), DNMT3A and HDAC2
show a subtype-specific co-expression coupling:

\begin{center}
\begin{tabular}{p{3cm}p{3cm}p{5.5cm}}
\toprule
\textbf{Subtype} &
\textbf{DNMT3A/HDAC2 coupling} &
\textbf{Biological meaning} \\
\midrule
Luminal~B &
$r = +0.267$ &
Active co-repressor circuit. DNMT3A and
HDAC2 are co-recruited to the same
chromatin targets. ER output is suppressed
by co-repressor complex activity. \\[0.5em]
Luminal~A &
$r = +0.071$ &
Co-repressor circuit absent. DNMT3A and
HDAC2 are not co-recruited. ER output
is not chromatically locked. \\
\midrule
\multicolumn{2}{l}{\textit{Difference}} &
$p = 5.68\times10^{-56}$ in METABRIC
($n=1{,}980$) \\
\bottomrule
\end{tabular}
\end{center}

\vspace{0.5em}

The coupling is $3.76\times$ stronger in LumB
than LumA. This is not a continuous gradation ---
it is a subtype-specific circuit that is switched
on in LumB and off in LumA.

\subsection{The Downstream Consequence:
TFF1/ESR1 Decoupling}

TFF1 is a canonical oestrogen-responsive gene:
when ER is active, TFF1 is transcribed. ESR1 is
the ER transcript itself.

In normal luminal biology: ESR1 $\uparrow$ implies
TFF1 $\uparrow$ (ER present and active).

In LumB: ESR1 is high (higher than LumA), but
TFF1 is specifically and significantly lower than
expected. The TFF1/ESR1 ratio is lower in LumB
than LumA despite LumB having \textit{higher}
raw ESR1. This is the chromatin lock working:
the ER is present and transcribed, but its
downstream output programme is suppressed.

\begin{center}
\colorbox{findingbox}{
\begin{minipage}{0.88\textwidth}
\vspace{0.5em}
\textbf{TFF1/ESR1 decoupling in METABRIC:}\\[0.3em]
$p=0.0019$, $n=1{,}980$\\[0.2em]
LumB: lower TFF1/ESR1 than LumA
\textit{despite} higher ESR1 transcript.\\[0.2em]
This is the measurable IHC fingerprint of the
DNMT3A/HDAC2 chromatin lock.\\[0.2em]
Full evidence: CS-LIT-9/23,
\href{https://doi.org/10.5281/zenodo.18884234}
{DOI: 10.5281/zenodo.18884234}
\vspace{0.5em}
\end{minipage}}
\end{center}

\subsection{Why Entinostat Specifically}

Entinostat is a class~I HDAC inhibitor (primarily
HDAC1, HDAC2, and HDAC3). It inhibits HDAC2
directly. The DNMT3A/HDAC2 co-repressor complex
requires HDAC2 enzymatic activity for its
chromatin-silencing function. When entinostat
inhibits HDAC2, the co-repressor complex loses
its functional capacity, HDAC2-dependent
acetylation marks are released, and the suppressed
ER output programme is restored.

\begin{center}
\colorbox{mechanismbox}{
\begin{minipage}{0.88\textwidth}
\vspace{0.5em}
\textbf{Mechanistic chain:}\\[0.4em]
DNMT3A/HDAC2 co-repressor suppresses
TFF1, GATA3, and ER target genes in LumB\\
$\downarrow$\\
Entinostat inhibits HDAC2 (class~I HDAC)\\
$\downarrow$\\
Co-repressor complex loses enzymatic function\\
$\downarrow$\\
ER output programme restored (TFF1, GATA3
re-express)\\
$\downarrow$\\
ER is now fully functional\\
$\downarrow$\\
Fulvestrant (or aromatase inhibitor) can
now engage the fully active receptor\\[0.4em]
\textbf{In LumA: the co-repressor circuit is
absent.} Entinostat has no HDAC2-dependent
chromatin lock to dissolve. The ER is already
fully functional. Entinostat adds no mechanism
for improving ET response in LumA.
\vspace{0.5em}
\end{minipage}}
\end{center}

% ══════════════════════════════════════════════════════════════
\section{The Published Evidence: What Is
Confirmed and What Is Novel}
% ══════════════════════════════════════════════════════════════

\subsection{What IS in the Literature}

\begin{enumerate}
    \item \textbf{China NMPA approval (April 2024):}
    Entinostat approved for HR$+$/HER2$-$ advanced
    breast cancer after prior endocrine therapy.
    First HDACi approval in breast cancer globally.

    \item \textbf{Meta-analysis (Springer, May 2025):}
    Pooled four RCTs ($n=1{,}371$). PFS HR$=0.80$,
    $p=0.01$ in HR$+$/HER2$-$ breast cancer. No OS
    benefit.

    \item \textbf{E2112 (SWOG/NRG):}
    Entinostat $+$ exemestane vs exemestane alone
    in advanced HR$+$ breast cancer post-AI. Negative
    for OS and PFS. Published. The meta-analysis
    aggregates E2112 and three additional RCTs to
    reach statistical significance.

    \item \textbf{DNMT3A/HDAC interactions
    (published biology):}
    DNMT3A and HDAC activity co-operate on
    co-repressor complexes to silence target genes
    (Narcancer 2021). HDAC2 is overexpressed in
    breast cancer and contributes to endocrine
    resistance (Springer 2024). DNMT3A
    misregulation is associated with tamoxifen
    resistance (OAE Publishing 2024).

    \item \textbf{NCT07235618 registered:}
    Phase~II entinostat $+$ fulvestrant, post-CDK4/6i
    failure, Sun Yat-sen University, start
    January 2026, $N=50$, primary endpoint PFS.
    PI: Dr.\ Shusen Wang (wangshs@sysucc.org.cn).
\end{enumerate}

\subsection{What is NOT in the Literature}

\begin{center}
\colorbox{novelbox}{
\begin{minipage}{0.88\textwidth}
\vspace{0.5em}
\textbf{The following have no published equivalent
as of 2026-03-06:}
\begin{itemize}[noitemsep]
    \item The quantification of the DNMT3A/HDAC2
    co-expression coupling difference between LumB
    ($r=+0.267$) and LumA ($r=+0.071$) at
    $p=5.68\times10^{-56}$.
    \item The framing of this coupling as a
    subtype-specific chromatin lock on ER output
    that explains the LumA/LumB paradox.
    \item The prediction that entinostat benefit
    is LumB-specific (not an HR$+$ class effect)
    because the DNMT3A/HDAC2 lock is present in
    LumB and absent in LumA.
    \item The proposal to stratify NCT07235618
    by LumA vs LumB to test this prediction.
    \item The use of TFF1/ESR1 ratio or
    FOXA1/EZH2 IHC ratio as the patient selection
    instrument for LumB enrichment.
    \item Any published trial stratifying entinostat
    benefit by LumA vs LumB PAM50 subtype.
\end{itemize}
\vspace{0.5em}
\end{minipage}}
\end{center}

% ══════════════════════════════════════════════════════════════
\section{Why the E2112 Failure and the Approval
Are Consistent}
% ══════════════════════════════════════════════════════════════

E2112 was negative for OS and PFS as a class-level
study (unselected HR$+$ population). The China NMPA
approval and the May 2025 meta-analysis show modest
but statistically significant PFS benefit
(HR$=0.80$) in pooled HR$+$ data.

The framework provides the reconciling explanation:
\begin{itemize}
    \item E2112 enrolled unselected HR$+$ patients ---
    approximately 50\% LumA and 50\% LumB (the
    typical PAM50 distribution in HR$+$ advanced
    breast cancer).
    \item LumA patients have no DNMT3A/HDAC2 lock.
    Entinostat provides no mechanism for improving
    their ET response. They are non-responders by
    mechanism.
    \item LumB patients have the DNMT3A/HDAC2 lock.
    Entinostat dissolves it. They are responders
    by mechanism.
    \item An approximately 50\%/50\% mixture of
    mechanism-responders (LumB) and
    non-responders (LumA) would produce an
    attenuated class-level effect ---
    consistent with HR$=0.80$.
\end{itemize}

The framework predicts that a LumB-enriched trial
would show HR significantly below $0.80$ ---
potentially approaching the magnitude seen in trials
of agents specifically matched to their mechanism,
such as palbociclib in LumA (HR$=0.58$, PALOMA-2).

% ══════════════════════════════════════════════════════════════
\section{Implications for NCT07235618}
% ══════════════════════════════════════════════════════════════

\begin{center}
\colorbox{actionbox}{
\begin{minipage}{0.88\textwidth}
\vspace{0.5em}
\textbf{NCT07235618 TRIAL DETAILS:}\\[0.4em]
\textbf{Title:} Phase~II study of entinostat
$+$ fulvestrant in HR$+$/HER2$-$ locally advanced
or metastatic breast cancer, recurrence or
progression after CDK4/6 inhibitor $+$ endocrine
therapy.\\[0.3em]
\textbf{Sponsor:} Sun Yat-sen University Cancer
Center, Guangzhou, China.\\[0.3em]
\textbf{PI:} Dr.\ Shusen Wang, MD\\[0.3em]
\textbf{PI contact:} wangshs@sysucc.org.cn\\[0.3em]
\textbf{Secondary contact:} Dr.\ Qiufan Zheng,
zhengqf@sysucc.org.cn\\[0.3em]
\textbf{Start date:} January 2026.
\textbf{$N=50$.}
\textbf{Primary endpoint:} PFS.\\[0.3em]
\textbf{Current biomarker:} PBMC histone
acetylation only.\\[0.4em]
\textbf{The LumB stratification hypothesis is
testable in this trial with:}
\begin{enumerate}[noitemsep]
    \item PAM50 subtype on baseline core biopsy
    (if tissue is banked)
    \item OR FOXA1/EZH2 IHC ratio on baseline
    tissue (two antibody stains; full protocol
    in CS-LIT-22,
    \href{https://doi.org/10.5281/zenodo.18892788}
    {DOI: 10.5281/zenodo.18892788})
    \item OR TFF1/ESR1 IHC ratio on baseline tissue
    (single additional stain; full protocol in
    CS-LIT-9/23,
    \href{https://doi.org/10.5281/zenodo.18884234}
    {DOI: 10.5281/zenodo.18884234})
\end{enumerate}
None of these require a protocol amendment,
change to primary endpoint, or additional
patient burden.
\vspace{0.5em}
\end{minipage}}
\end{center}

\subsection{Why the Post-CDK4/6i Population
Is De Facto LumB-Enriched}

CDK4/6 inhibitor $+$ endocrine therapy is
highly effective in LumA (CDKN1A-mediated cell
cycle arrest; the CDK4/6 lock is the primary
mechanism in LumA). Patients progressing on
CDK4/6i $+$ ET are disproportionately LumB:
LumA patients with intact CDKN1A response are
more likely to respond and remain on therapy.
LumB patients with the DNMT3A/HDAC2 lock
suppressing ER output have a mechanism of
resistance that CDK4/6i does not address, making
progression more likely.

NCT07235618's eligibility criterion (prior CDK4/6i
failure) therefore selects a population enriched
for LumB biology --- the exact subtype the
framework predicts will benefit from entinostat.
The trial is inadvertently well-designed for
the hypothesis, even without subtype stratification
built in.

% ══════════════════════════════════════════════════════════════
\section{Validation Pathway}
% ══════════════════════════════════════════════════════════════

\begin{center}
\colorbox{actionbox}{
\begin{minipage}{0.88\textwidth}
\vspace{0.5em}
\textbf{Proposed analyses:}\\[0.4em]
\textbf{Analysis 1 --- NCT07235618 correlative:}\\
Add PAM50 or FOXA1/EZH2 IHC ratio to baseline
tissue. Compare PFS in LumB vs LumA subgroups.
Falsifiable prediction: LumB-stratified HR
$<0.80$; LumA-stratified HR $\approx 1.0$.\\[0.4em]
\textbf{Analysis 2 --- E2112 post-hoc subgroup:}\\
Re-analyse E2112 dataset with PAM50 stratification
(if tissue banks are accessible). LumB subgroup
response should be substantially larger than the
class-level HR$=0.80$.\\[0.4em]
\textbf{Analysis 3 --- TFF1/ESR1 as continuous
predictor:}\\
In NCT07235618 or E2112 tissue: test TFF1/ESR1
IHC ratio as a continuous predictor of entinostat
benefit magnitude. Lower ratio (deeper chromatin
lock) should predict greater benefit.\\[0.4em]
\textbf{Analysis 4 --- DNMT3A/HDAC2 IHC:}\\
Develop DNMT3A/HDAC2 co-expression IHC as a
direct biomarker for entinostat eligibility.
Correlation with TFF1/ESR1 ratio and PAM50 LumB
classification would establish the mechanistic
chain at the tissue level.
\vspace{0.5em}
\end{minipage}}
\end{center}

% ══════════════════════════════════════════════════════════════
\section{What Is and Is Not Claimed}
% ══════════════════════════════════════════════════════════════

\subsection{What IS Claimed}

\begin{itemize}
    \item DNMT3A/HDAC2 co-expression coupling is
    $3.76\times$ stronger in LumB than LumA
    ($r=+0.267$ vs $r=+0.071$,
    $p=5.68\times10^{-56}$, METABRIC $n=1{,}980$).
    \item This coupling is the mechanistic basis of
    ER output suppression in LumB (TFF1/ESR1
    decoupling, $p=0.0019$).
    \item Entinostat (class~I HDACi) directly
    inhibits HDAC2 and is mechanistically predicted
    to dissolve this lock specifically in LumB.
    \item The class-level HR$=0.80$ from published
    meta-analysis is consistent with a strong LumB
    signal diluted by non-responding LumA patients.
    \item NCT07235618 is a de facto LumB-enriched
    population and is the most immediate trial in
    which this prediction is testable.
    \item No published trial has stratified
    entinostat benefit by LumA vs LumB.
\end{itemize}

\subsection{What is NOT Claimed}

\begin{itemize}
    \item The LumB-specific benefit has not been
    prospectively confirmed in any clinical trial.
    \item The DNMT3A/HDAC2 coupling has not been
    validated as a predictive biomarker on stained
    tissue.
    \item This does not replace or supersede the
    existing entinostat literature; it extends
    it with subtype-specific mechanistic framing.
    \item Clinical implementation requires
    prospective subtype-stratified validation.
\end{itemize}

% ══════════════════════���═══════════════════════════════════════
\section{Locked Statement}
% ══════════════════════════════════════════════════════════════

\begin{center}
\colorbox{findingbox}{
\begin{minipage}{0.88\textwidth}
\vspace{0.6em}
\textbf{Stated once, precisely, without
inflation:}\\[0.5em]
Waddington attractor geometry identifies a
LumB-specific DNMT3A/HDAC2 chromatin co-repressor
circuit ($r=+0.267$ LumB vs $r=+0.071$ LumA,
$p=5.68\times10^{-56}$, METABRIC $n=1{,}980$)
that suppresses ER output in LumB despite elevated
ESR1 transcript (TFF1/ESR1 decoupling, $p=0.0019$).
This circuit is absent in LumA. Entinostat, a
class~I HDAC inhibitor (HDAC2 target), is
mechanistically predicted to dissolve this lock
specifically in LumB --- producing an entinostat
benefit that is LumB-enriched, not an HR$+$ class
effect. The class-level PFS HR$=0.80$ from the
published meta-analysis ($n=1{,}371$) is consistent
with a strong LumB signal diluted by non-responding
LumA patients. NCT07235618 (entinostat $+$
fulvestrant, post-CDK4/6i, Sun Yat-sen, $N=50$,
start January 2026) is a de facto LumB-enriched
population in which this prediction is immediately
testable as a correlative endpoint without protocol
amendment. No published trial has stratified
entinostat benefit by LumA vs LumB. The subtype-specific
mechanistic prediction has no direct published
equivalent as of 2026-03-06.
\vspace{0.5em}
\end{minipage}}
\end{center}

% ══════════════════════════════════════════════════════════════
\section{Document Metadata}
% ═════════════════════════���════════════════════════════════════

\begin{tabular}{ll}
\toprule
\textbf{Field} & \textbf{Value} \\
\midrule
Document ID & CS-LIT-15 \\
Parent document & BRCA-S8h (2026-03-05) \\
Verdict & CONVERGENT-NOVEL \\
Date & 2026-03-06 \\
Author & Eric Robert Lawson \\
ORCID & 0009-0002-0414-6544 \\
Contact & OrganismCore@proton.me \\
Repository &
\href{https://github.com/Eric-Robert-Lawson/attractor-oncology}
{attractor-oncology} \\
Derivation data &
GSE176078 (19,542 single cancer cells) \\
Replication data &
METABRIC ($n=1{,}980$) \\
Key statistics &
$r=+0.267$ LumB vs $r=+0.071$ LumA, \\
& $p=5.68\times10^{-56}$;
TFF1/ESR1 $p=0.0019$ \\
Trial &
NCT07235618 (Sun Yat-sen, start Jan 2026) \\
Trial PI &
Dr.\ Shusen Wang
(wangshs@sysucc.org.cn) \\
CS-LIT-9/23 DOI &
\href{https://doi.org/10.5281/zenodo.18884234}
{10.5281/zenodo.18884234} \\
CS-LIT-1 DOI &
\href{https://doi.org/10.5281/zenodo.18883922}
{10.5281/zenodo.18883922} \\
CS-LIT-2 DOI &
\href{https://doi.org/10.5281/zenodo.18884158}
{10.5281/zenodo.18884158} \\
CS-LIT-22 DOI &
\href{https://doi.org/10.5281/zenodo.18892788}
{10.5281/zenodo.18892788} \\
License & CC BY 4.0 \\
\bottomrule
\end{tabular}

\end{document}
